% =============================================================================
%  CLAS12 Technical Note --- references
%
%  The recurring foundational CLAS12-note references reused across the beam-
%  background notes (authoring guide, section 7). Cite with \cite{key}; keep the
%  entries you use and add note-specific references below.
% =============================================================================

\begin{thebibliography}{9}

\bibitem{note:moller}
  R.~De Vita and M.~Ungaro,
  {\it CLAS12-note 2016-006: Moller shield --- GEMC-optimized layout vs.\ the
  engineering design.}

\bibitem{note:beamline}
  R.~De Vita {\it et al.},
  {\it CLAS12-note 2017-012: Corrections to the CLAS12 vacuum beamline.}

\bibitem{note:target}
  M.~Ungaro,
  {\it CLAS12-note 2017-013: Corrections to the CLAS12 target design.}

\bibitem{note:embg}
  R.~De Vita {\it et al.},
  {\it CLAS12-note 2017-016: Electromagnetic background rates in CLAS12.}

\bibitem{note:cad}
  R.~De Vita and M.~Ungaro,
  {\it CLAS12-note 2017-017: Importing the CLAS12 CAD models of the target and
  beamline.}

\bibitem{note:ftof}
  D.~Carman,
  {\it FTOF-note 2014-06: Forward Time-of-Flight PMT currents.}

\bibitem{gemc:epj}
  M.~Ungaro,
  {\it GEMC: A database-driven simulation program,}
  EPJ Web Conf.\ {\bf 295} (2024).

\bibitem{gemc:nim}
  M.~Ungaro {\it et al.},
  {\it The CLAS12 Geant4 simulation,}
  Nucl.\ Instrum.\ Meth.\ A {\bf 959} (2020) 163422.

\end{thebibliography}
